% !TeX spellcheck = en_US

\ClearPage
\pdfbookmark[1]{Abstract}{sec:abstract-en}
\chapter*{Abstract}
In machine learning applications for climate research, the selection and
understanding of input features is a critical task for effective model
development. This is particularly challenging for meteorological reanalysis datasets such as ERA5, which offer a large number of atmospheric and
land-surface variables at high spatial and temporal resolution, resulting in data volumes that prohibit undirected exploration (see Chapter~\ref{ch:introduction}).

This thesis presents \textit{glaciVis}, an interactive visualization tool designed to support structured feature exploration of ERA5 data in the context of glacier research. The tool combines a preprocessing pipeline for converting ERA5 and WGMS glacier data into analysis-ready formats with a coordinated spatial-temporal interface built on Panel and Bokeh.
Users can select datasets, configure filtering and aggregation parameters, inspect spatial patterns on an interactive map, and analyze temporal trends through linked time series plots, including overlay of glacier mass balance
observations.

The system is demonstrated through a use-case investigating the relationship between wind patterns and glacier mass balance in the Alps.
While the explored hypothesis did not reveal a clear wind-driven signal, the investigation illustrated how the tool enables a systematic and visual hypothesis testing before going to feature selection for predictive modeling.
The resulting framework provides a foundation for extending the
analysis to additional ERA5 variables and integrating automated statistical methods in future work.